\chapter{Introduction}

roadrunner\_nitro is an engine for simulating biochemical models written in SBML, the Systems Biology Markup
Language. It does the same job as libroadrunner, gives the same answers, and is written in Rust. It needs neither
libroadrunner nor libsbml: it reads SBML itself and compiles each model to native code when it loads.

\section{What it does}

\begin{itemize}
  \item \textbf{Time courses}: deterministic simulation with stiff and non-stiff solvers, events, delays,
        assignment and rate rules; stochastic simulation with Gillespie's method.
  \item \textbf{Steady states}: Newton's method on the system reduced by its conservation laws, with integration
        towards the steady state when Newton's method alone fails.
  \item \textbf{Metabolic control analysis}: Jacobians, elasticities, flux and concentration control
        coefficients.
  \item \textbf{Parameter estimation}: fitting models to data with exact gradients, local and global optimizers,
        and confidence intervals; problems can be built in Python or C, or read from PEtab files.
\end{itemize}

\section{Two ways to use it}

\begin{center}
\begin{tabularx}{\textwidth}{L{0.33\textwidth}Y}
\toprule
If you write \ldots & Use \\
\midrule
Python & the Python package \code{roadrunner\_nitro} (chapters~\ref{ch:simulation} and~\ref{ch:fitting}) \\
\addlinespace
C, C++, MATLAB, Julia, C\#, R, or anything that can load a C library &
  the C library \code{libroadrunner\_nitro\_c\_api}, which has all 350 functions of libroadrunner's C API
  (chapter~\ref{ch:capi}) \\
\bottomrule
\end{tabularx}
\end{center}

The Python methods have libroadrunner's names (\code{simulate}, \code{steadyState}, \code{getCC},
\code{setValue}, \ldots), and the C library has libroadrunner's C functions with the same names, arguments and
memory rules. Programs written for libroadrunner therefore run on roadrunner\_nitro with few or no changes;
chapter~\ref{ch:differences} lists the places where it deliberately behaves differently.

roadrunner\_nitro is the engine inside Tellurium Nitro. If you want the whole Tellurium environment in Python
(\code{te.loada}, plotting, SED-ML, COMBINE archives), install Tellurium Nitro, which includes this engine.

The package has its own name, \code{roadrunner\_nitro}, so an existing libroadrunner installation is not
touched and both can be used side by side.

\section{About this manual}

Chapter~\ref{ch:installation} installs the Python package and the C library. Chapter~\ref{ch:simulation} covers
simulation, steady states and control analysis from Python. Chapter~\ref{ch:fitting} covers parameter
estimation, from a three-line fit of one time course to PEtab problems with many experiments.
Chapter~\ref{ch:capi} shows the C library, and chapter~\ref{ch:differences} the differences from libroadrunner.

Outputs shown in this manual were produced by the code beside them. Complete programs are in the repository's
\code{examples/python} and \code{capi/examples} folders.
